\documentclass{colm2026}
% For blind review: \documentclass[review]{colm2026}

\usepackage{times}
\usepackage{latexsym}
\usepackage[T1]{fontenc}
\usepackage[utf8]{inputenc}
\usepackage{microtype}
\usepackage{inconsolata}
\usepackage{booktabs}
\usepackage{graphicx}
\usepackage{amsmath}
\usepackage{amssymb}
\usepackage{multirow}
\usepackage{subcaption}
\usepackage{hyperref}
\usepackage{xcolor}

\title{Do VLMs Actually Look? \\
Grounding Faithfulness Across Reasoning Steps in Vision-Language Models}

\author{Anonymous}

\begin{document}
\maketitle

% ── Abstract ──────────────────────────────────────────────────────────────────
\begin{abstract}
[TODO: 150--200 words.
State the problem: VLMs may generate chain-of-thought reasoning that drifts from
visual evidence as step depth increases.
State the method: we introduce the Grounding Faithfulness Score (GFS), a step-wise
metric computed from vision-encoder attribution maps (Grad-CAM, attention rollout,
patch occlusion) and validated by image-swap sensitivity.
State key results: GFS decays significantly with step index across all three models
(mean slope = [X], $p < 0.05$); swap sensitivity confirms that [N]\% of steps
beyond step 3 are visually ungrounded.
State implication: multimodal CoT elicits surface reasoning chains, not
visually-grounded propositional inference.]
\end{abstract}

% ── 1. Introduction ───────────────────────────────────────────────────────────
\section{Introduction}
\label{sec:intro}

[TODO: Motivate the question. VLMs prompted for CoT produce multi-step reasoning,
but it is unclear whether later steps are grounded in the image or are purely
language-driven confabulation.

Key gap: existing hallucination benchmarks (POPE, CHAIR) measure final-answer
accuracy, not step-level grounding fidelity.

Contributions bullet list:
\begin{itemize}
  \item We introduce \textbf{GFS}, a step-wise grounding faithfulness metric
        derived from vision-encoder attribution maps.
  \item We construct a 200-triplet swap-pair benchmark from ScienceQA and GQA
        that validates GFS via controlled image substitution.
  \item We run GFS experiments on three 7--8B VLMs (LLaVA-1.5, InternVL2, Idefics2)
        and show a significant decay pattern with step depth.
  \item We release code, benchmark, and results at [URL].
\end{itemize}]

% ── 2. Related Work ───────────────────────────────────────────────────────────
\section{Related Work}
\label{sec:related}

\paragraph{VLM hallucination.}
[TODO: Cite POPE \citep{li2023pope}, object hallucination \citep{rohrbach2018hallucination},
Favero et al. \citep{favero2024multimodal}. Distinguish: prior work measures
output-level hallucination; we measure step-level attribution faithfulness.]

\paragraph{Multimodal chain-of-thought.}
[TODO: Cite Wei et al. \citep{wei2022chain}, Zhang et al. \citep{zhang2023multimodal_cot},
LLaVA-CoT \citep{bai2024llava_cot}. Gap: none measure whether steps are
visually grounded step-by-step.]

\paragraph{Attention and gradient attribution in transformers.}
[TODO: Cite Grad-CAM \citep{selvaraju2017gradcam}, rollout \citep{abnar2020rollout},
Parcalabescu \citep{parcalabescu2023measuring}. Our work applies these to
multi-step VLM inference.]

\paragraph{Faithfulness of explanations.}
[TODO: Cite Turpin et al. \citep{agrawal2024faithfulness}.
Position: their work is NLP-only; we extend to vision encoders.]

% ── 3. Method ─────────────────────────────────────────────────────────────────
\section{Method}
\label{sec:method}

\subsection{Task Formulation}
\label{sec:formulation}

[TODO: Define formally. Input: image $I$, question $Q$.
VLM generates CoT: $C = (s_1, s_2, \ldots, s_K, a)$ where $s_k$ is step $k$
and $a$ is the final answer.
Goal: measure visual grounding of each step $s_k$.]

\subsection{Grounding Faithfulness Score (GFS)}
\label{sec:gfs}

For each step $s_k$, we compute an attribution heatmap $H_k \in \mathbb{R}^{W \times W}$
from the vision encoder and define:
\begin{equation}
  \text{GFS}_k = 1 - \frac{-\sum_{i,j} p_{ij} \log p_{ij}}{\log(W^2)}
  \label{eq:gfs}
\end{equation}
where $p_{ij} = H_k[i,j] / \sum_{i',j'} H_k[i',j']$.
High GFS indicates concentrated, spatially specific attention.
Low GFS indicates diffuse attention, consistent with language-driven generation.

Steps without explicit visual references (detected via noun-phrase analysis) are
excluded from GFS computation and marked as \texttt{None}.

\subsection{Attribution Methods}
\label{sec:attribution}

We compute $H_k$ using three independent methods:

\paragraph{Grad-CAM.} [TODO: Describe implementation: hook last vision encoder layer,
backpropagate from final-token logit sum, weight activations by gradients.]

\paragraph{Attention Rollout.} [TODO: Cite \citep{abnar2020rollout}. Describe:
chain attention matrices across all LM layers with residual addition.]

\paragraph{Patch Occlusion.} [TODO: Describe: $24 \times 24$ patches, stride 12,
score = baseline logit $-$ occluded logit.]

\subsection{Swap-Pair Benchmark}
\label{sec:benchmark}

[TODO: Describe triplet construction. Same subject, different answer, pixel
MSE $>$ $(0.25)^2$ on $224 \times 224$ resize. 200 triplets from ScienceQA
(70\%) and GQA (30\%). Define swap sensitivity score.]

\begin{table}[t]
\centering
\caption{Benchmark composition.}
\label{tab:benchmark}
\begin{tabular}{lcc}
\toprule
Source & Triplets & Subjects \\
\midrule
ScienceQA & [TODO] & [TODO] \\
GQA & [TODO] & [TODO] \\
\midrule
Total & 200 & [TODO] \\
\bottomrule
\end{tabular}
\end{table}

% ── 4. Experimental Setup ─────────────────────────────────────────────────────
\section{Experimental Setup}
\label{sec:setup}

\subsection{Models}
[TODO: LLaVA-1.5-7B \citep{liu2024llava}, InternVL2-8B \citep{chen2024internvl2},
Idefics2-8B \citep{laurenccon2024idefics2}. All loaded in 4-bit NF4 quantization
on a single A100 40GB GPU. Greedy decoding, \texttt{max\_new\_tokens=512}.]

\subsection{CoT Prompting}
[TODO: Template from Section~\ref{sec:formulation}. Force ``Step 1:'' prefix.
Parse steps with regex + fallback on double-newline split.]

\subsection{Evaluation Metrics}
[TODO: (1) Mean GFS per step, (2) decay slope (linear fit over steps), (3) swap
sensitivity score, (4) final-answer accuracy vs ground truth.]

% ── 5. Results ────────────────────────────────────────────────────────────────
\section{Results}
\label{sec:results}

\subsection{GFS Decays With Step Depth (RQ1)}

[TODO: Report decay slope, CI, t-test results. Reference Figure~\ref{fig:gfs_decay}.]

\begin{figure}[t]
\centering
\includegraphics[width=0.9\linewidth]{../results/figures/fig1_gfs_decay.pdf}
\caption{Mean GFS per reasoning step across three models (shaded = $\pm$1 SE).
GFS decays significantly with step depth, indicating reduced visual grounding
in later reasoning steps.}
\label{fig:gfs_decay}
\end{figure}

\subsection{Swap Sensitivity Confirms Ungrounded Steps (RQ2)}

[TODO: Report swap sensitivity scores. Reference Figure~\ref{fig:swap}.]

\begin{figure}[t]
\centering
\includegraphics[width=0.9\linewidth]{../results/figures/fig2_swap_sensitivity.pdf}
\caption{GFS vs swap sensitivity score per sample. High swap sensitivity =
model actively used the image. Low sensitivity = model ignored the image
after early steps.}
\label{fig:swap}
\end{figure}

\subsection{Main Results Table}

\begin{table*}[t]
\centering
\caption{Main results: GFS, decay slope, swap sensitivity, and accuracy.
* = $p < 0.05$ on one-sample t-test of per-sample decay slopes against 0.}
\label{tab:main}
\begin{tabular}{lccccc}
\toprule
Model & Mean GFS (95\% CI) & Decay slope & Swap sens. & Accuracy & $p$ \\
\midrule
LLaVA-1.5-7B   & [TODO] & [TODO] & [TODO] & [TODO] & [TODO] \\
InternVL2-8B   & [TODO] & [TODO] & [TODO] & [TODO] & [TODO] \\
Idefics2-8B    & [TODO] & [TODO] & [TODO] & [TODO] & [TODO] \\
\bottomrule
\end{tabular}
\end{table*}

\subsection{Heatmap Qualitative Analysis}

[TODO: Reference Figure~\ref{fig:heatmaps}. Describe pattern: step 1 heatmaps
concentrate on task-relevant regions; step 4+ heatmaps are diffuse.]

\begin{figure*}[t]
\centering
\includegraphics[width=0.95\linewidth]{../results/figures/fig3_heatmap_grid.pdf}
\caption{Attribution heatmaps per reasoning step for two representative samples
(rows). GFS values annotated per step. Visual grounding degrades with step depth.}
\label{fig:heatmaps}
\end{figure*}

% ── 6. Analysis ───────────────────────────────────────────────────────────────
\section{Analysis}
\label{sec:analysis}

\subsection{Ablation A1: CoT vs Direct-Answer Prompting}

[TODO: Report mean GFS for CoT vs direct. Does forcing CoT help or hurt grounding?
Reference Figure~\ref{fig:cot_vs_direct}.]

\begin{figure}[t]
\centering
\includegraphics[width=0.8\linewidth]{../results/figures/fig4_cot_vs_direct.pdf}
\caption{GFS distribution for CoT vs direct-answer prompting (LLaVA-1.5-7B).}
\label{fig:cot_vs_direct}
\end{figure}

\subsection{Ablation A2: Attribution Method Comparison}

[TODO: Spearman correlation between GFS from Grad-CAM vs Rollout vs Occlusion.
Do methods agree? Robustness check.]

\begin{table}[t]
\centering
\caption{Spearman $\rho$ between attribution methods on mean GFS (N=30 samples).}
\label{tab:methods}
\begin{tabular}{lcc}
\toprule
 & Rollout & Occlusion \\
\midrule
Grad-CAM & [TODO] & [TODO] \\
Rollout  & —      & [TODO] \\
\bottomrule
\end{tabular}
\end{table}

\subsection{Ablation A3: GFS by Question Category}

[TODO: Table or figure showing GFS by ScienceQA subject (biology, physics, chemistry).
Does category moderate the decay?]

% ── 7. Discussion ─────────────────────────────────────────────────────────────
\section{Discussion}
\label{sec:discussion}

[TODO: Interpret findings.
Why does grounding decay? Hypotheses: (1) LM component dominates after step 1--2
due to autoregressive inertia; (2) image tokens are down-weighted in later layers;
(3) CoT prompting encourages language-pattern completion, not visual inspection.

Implications: VLM CoT is not a reliable proxy for visual reasoning.
Design recommendation: re-attend to image at each step (e.g. image token refresh).

Connect to faithfulness literature: parallels Turpin et al. unfaithful NLP CoT,
but in multimodal setting.]

% ── 8. Conclusion ─────────────────────────────────────────────────────────────
\section{Conclusion}
\label{sec:conclusion}

[TODO: Summarize. GFS decay is consistent, significant, model-agnostic.
Swap sensitivity confirms later steps are not image-conditioned.
Future work: (1) mechanistic analysis of image-token attention in late layers;
(2) training intervention to enforce per-step grounding; (3) extension to video VLMs.]

% ── Limitations ───────────────────────────────────────────────────────────────
\section{Limitations}
\label{sec:limitations}

[TODO:
\begin{itemize}
  \item GFS is a proxy metric (entropy of attribution) not a direct measure
        of semantic alignment between step text and image region.
  \item Attribution methods may be unreliable on 4-bit quantized models.
  \item 200 samples is a modest scale; effect sizes may shift with larger benchmarks.
  \item We study three English-only VLMs; generalization to multilingual or
        domain-specific VLMs is unknown.
\end{itemize}]

% ── Ethics ────────────────────────────────────────────────────────────────────
\section{Ethics Statement}
\label{sec:ethics}

[TODO: All datasets (ScienceQA, GQA) are publicly available under standard
research licenses. No personal data. No crowdsourcing. Results reveal a
limitation of VLMs that could inform responsible deployment in high-stakes
visual reasoning tasks (medical, scientific). Code will be released under MIT license.]

% ── References ────────────────────────────────────────────────────────────────
\bibliographystyle{colm2026_natbib}
\bibliography{references}

\end{document}
